\begin{enumerate}
%1
\item \textbf{ทบทวนปริพันธ์} จงหาปริพันธ์ต่อไปนี้

\begin{tasks}[after-skip = 80pt, after-item-skip = 80pt](3)
\task $\int 1 \, dx$
\task $\int x^2 \, dx$
\task $\int x^{3/4} \, dx$
\task $\int \frac{1}{x^3} \, dx$
\task $\int (12x^4 + 12x^3) \, dx$
\task $\int 4e^x \,dx$
\task $\int 3\sin x \,dx$
\task $\int -2\cos x \,dx$
\task $\int \frac{2}{x} \,dx$
\end{tasks}

\item \textbf{ค่าเชิงอนุพันธ์} ถ้ากำหนดฟังก์ชัน $u(x)$ ที่ทำให้ $\frac{du}{dx} = u'(x)$ แล้วค่าเชิงอนุพันธ์คือ
\[ du = \rule{100pt}{0.5pt}\]

จงหาค่าเชิงอนุพันธ์ของฟังก์ชันต่อไปนี้

\begin{tasks}[after-skip = 60pt, after-item-skip=60pt](3)
\task $u(x) = 3x+2$
\task $u(x) = x^3+x^2+1$
\task $u(x) = \sqrt{x}$
\task $u(x) = \frac{2}{x^3}$
\task $u(x) = \sin(2x)$
\task $u(x) = e^{3x+1}$
\end{tasks}

\newpage 
\item  \textbf{เทคนิคการแทนค่า $u$}

จงหาปริพันธ์ของฟังก์ชันประกอบต่อไปนี้ ด้วยการแทนค่า $u$
\begin{tasks}[after-skip = 100pt, after-item-skip=100pt](2)
\task $\int (2x+1)^4 \, dx$
\task $\int (1-3x)^6 \, dx$
\task $\int \sqrt{5x+2} \, dx$
\task $\int \frac{1}{(4x-2)^7} \, dx$
\task $\int x(x^2+1)^6 \, dx$
\task $\int \frac{x^2}{(x^3+1)^4} \, dx$
\task $\int \cos x (\sin x)^6 \, dx$
\task $\int xe^{x^2} \, dx$
\end{tasks}

\item \textbf{ปริพันธ์ฟังก์ชันอดิศัย}

จงใช้สูตรปริพันธ์ฟังก์ชันอดิศัย และการแทนค่า $u$ เพื่อหาปริพันธ์ต่อไปนี้

\begin{tasks}[after-skip = 60pt, after-item-skip=60pt](4)
\task $\int e^{4x} \, dx$
\task $\int e^{-3x-3} \, dx$
\task $\int \frac{2}{3x+4} \, dx$
\task $\int \frac{x}{1+x^2} \, dx$
\task $\int \sin (2x) \, dx$
\task $\int \cos (5x) \, dx$
\task $\int \sin (3x-1) \, dx$
\task $\int 4\cos \left( \frac{x+5}{2} \right) \, dx$
\end{tasks}


\end{enumerate}